% !TEX program = pdflatex
\documentclass[12pt,a4paper]{report}

% ============================================================
% PREAMBLE
% ============================================================
\usepackage[utf8]{inputenc}
\usepackage[T1]{fontenc}
\usepackage[margin=2.5cm]{geometry}
\usepackage{amsmath,amssymb,amsthm}
\usepackage{mathtools}
\usepackage{hyperref}
\usepackage{xcolor}
\usepackage{graphicx}
\usepackage{float}
\usepackage{booktabs}
\usepackage{enumitem}
\usepackage{microtype}
\usepackage{parskip}
\usepackage{titlesec}
\usepackage{fancyhdr}

% Listings for code
\usepackage{listings}
\lstset{
  basicstyle=\ttfamily\small,
  keywordstyle=\color{blue!70!black}\bfseries,
  commentstyle=\color{gray},
  stringstyle=\color{orange!80!black},
  numbers=left,
  numberstyle=\tiny\color{gray},
  numbersep=8pt,
  breaklines=true,
  showstringspaces=false,
  frame=single,
  rulecolor=\color{black!20},
  backgroundcolor=\color{gray!5},
  captionpos=b,
  tabsize=4,
  language=Python,
}

% tcolorbox for definitions and key insights
\usepackage{tcolorbox}
\tcbuselibrary{theorems,skins,breakable}

\newtcolorbox{definition}[2][]{
  colback=blue!5!white,
  colframe=blue!60!black,
  fonttitle=\bfseries,
  title={Definition: #2},
  #1,
  breakable
}

\newtcolorbox{insight}[2][]{
  colback=green!5!white,
  colframe=green!50!black,
  fonttitle=\bfseries,
  title={Key Insight: #2},
  #1,
  breakable
}

\newtcolorbox{codebox}[2][]{
  colback=gray!5!white,
  colframe=black!40,
  fonttitle=\bfseries,
  title={#2},
  #1,
  breakable
}

% Algorithm
\usepackage[ruled,vlined,linesnumbered]{algorithm2e}

% TikZ for diagrams
\usepackage{tikz}
\usetikzlibrary{arrows.meta,positioning,shapes,fit,backgrounds,calc}

% Math operators
\DeclareMathOperator*{\argmax}{arg\,max}
\DeclareMathOperator*{\argmin}{arg\,min}
\newcommand{\R}{\mathbb{R}}
\newcommand{\E}{\mathbb{E}}
\newcommand{\bh}{\mathbf{h}}
\newcommand{\bx}{\mathbf{x}}
\newcommand{\bg}{\mathbf{g}}
\newcommand{\bW}{\mathbf{W}}
\newcommand{\calV}{\mathcal{V}}
\newcommand{\calE}{\mathcal{E}}
\newcommand{\calG}{\mathcal{G}}
\newcommand{\calL}{\mathcal{L}}
\newcommand{\calN}{\mathcal{N}}
\newcommand{\calB}{\mathcal{B}}

% Headers
\pagestyle{fancy}
\fancyhf{}
\fancyhead[L]{\leftmark}
\fancyhead[R]{\thepage}
\renewcommand{\headrulewidth}{0.4pt}

% ============================================================
\begin{document}
% ============================================================

\begin{titlepage}
  \centering
  \vspace*{2cm}
  {\LARGE\bfseries Technical Guide\par}
  \vspace{0.5cm}
  {\Large GNN-Based Epidemic Source Detection\\on Temporal Contact Networks\par}
  \vspace{1.5cm}
  {\large Master's Thesis — Implementation Reference\par}
  \vspace{0.5cm}
  {\large Author: [Your Name]\par}
  \vspace{0.5cm}
  {\large\today\par}
  \vspace{2cm}
  \begin{abstract}
    \noindent
    This document is a self-contained technical reference for the master's thesis
    project on graph neural network (GNN) based epidemic source detection on
    temporal contact networks. It covers every layer of the software stack: the
    stochastic SIR epidemic model and its C-based simulation engine, the five
    distinct graph representations derived from a temporal contact sequence, the
    five GNN architectures (StaticGNN, TemporalGNN, BacktrackingNetwork, DBGNN,
    and DAG-GNN), the training pipeline with its loss function and early-stopping
    scheme, the evaluation metrics, the non-ML baselines, and the experimental
    design including the two parameter sweeps. Mathematical derivations are given
    in full, code listings are drawn directly from the implementation, and the
    rationale behind every design decision is explained.
  \end{abstract}
  \vfill
\end{titlepage}

\tableofcontents
\clearpage

% ============================================================
\chapter{Introduction and Problem Formulation}
\label{chap:intro}
% ============================================================

\section{Motivation}

Epidemic source detection is the problem of identifying, after an outbreak has
partially or fully unfolded, which individual node in a contact network most
likely initiated the spreading process. The practical stakes are high: in a
hospital ward an unidentified index patient can sustain a nosocomial infection
chain; in an organisation an undetected patient zero prolongs quarantine
decisions; in digital contagion (misinformation, computer malware) source
identification enables targeted countermeasures. Classical network epidemiology
provides analytical upper bounds and centrality heuristics, but these are
derived under strong structural assumptions (trees, homogeneous mixing,
infinite-graph limits) that break down on real finite contact networks.

Graph neural networks offer a data-driven alternative: by training on many
simulated outbreaks one can learn a discriminative function that maps an
observed SIR snapshot directly to a probability distribution over candidate
sources without assuming any particular network structure.

\section{The Temporal Dimension}

Standard GNN-based source detection approaches operate on the \emph{static
projection} of the contact network, aggregating all contacts across time into a
single undirected graph. This discards the most informative signal: the
\emph{ordering} of contacts. On a temporal contact network the set of causally
possible transmission paths is strictly smaller than the set of paths in the
static projection, because a chain $u \to v \to w$ can only transmit an
infection if the contact $(u,v)$ precedes the contact $(v,w)$ in real time. An
architecture that is unaware of this constraint must learn to ignore spurious
paths from data alone, which is both data-inefficient and prone to
generalisation failures when the temporal ordering changes.

Three of the five architectures in this project — BacktrackingNetwork,
DBGNN, and DAG-GNN — incorporate temporal ordering explicitly, either through
edge feature encoding (BacktrackingNetwork), a higher-order De Bruijn graph
(DBGNN), or a temporal event directed acyclic graph (DAG-GNN).

\section{Formal Problem Statement}

Let $G = (\calV, \calE^t)$ be a temporal contact network where $\calV =
\{1,\ldots,N\}$ is a fixed node set and $\calE^t = \{(u,v,t)\}$ is a
multi-set of timestamped undirected contact events. We assume discrete time
$t \in \{0,1,\ldots,T\}$.

\begin{definition}{Source Detection Problem}
Given:
\begin{itemize}[noitemsep]
  \item The temporal contact network $G = (\calV, \calE^t)$,
  \item An observed compartmental state $\mathbf{O} = (\mathbf{S}_{\text{obs}},
    \mathbf{I}_{\text{obs}}, \mathbf{R}_{\text{obs}}) \in \{0,1\}^{N \times 3}$
    recorded at time $t_{\text{end}}$, where $O_{v,k}=1$ iff node $v$ is in
    compartment $k \in \{\mathrm{S},\mathrm{I},\mathrm{R}\}$,
\end{itemize}
Find the most probable source node:
\begin{equation}
  s^* = \argmax_{v \in \calV} \; P(\text{source} = v \mid \mathbf{O}, G).
  \label{eq:source_detection}
\end{equation}
\end{definition}

Because computing the exact posterior in~\eqref{eq:source_detection} is
\#P-hard for general networks (requiring marginalisation over all stochastic
trajectories), we instead train a discriminative neural network
$f_\theta: \mathbf{O} \mapsto \log P_\theta(\text{source} = v \mid \mathbf{O})$
using synthetic ground-truth labels from Monte Carlo (MC) SIR simulations.

\section{Pipeline Overview}

The full codebase implements a three-stage pipeline:

\begin{enumerate}
  \item \textbf{Simulation} (\texttt{main\_tsir.py}): Run the C-based temporal
    SIR simulator for every possible source node $v \in \calV$ with $n_{\text{runs}}$
    stochastic realisations (ground-truth data) and $n_{\text{mc}}$ MC realisations
    (training data). Results are stored as binary arrays and logged as a versioned
    \texttt{wandb} artefact.
  \item \textbf{Training and Inference} (\texttt{main\_train.py}): Load the
    artefact, build the appropriate graph representation, train a GNN model on
    the MC simulations (cross-entropy loss), then run inference on the
    ground-truth simulations to produce per-run source probability vectors.
  \item \textbf{Evaluation} (\texttt{main\_eval.py}): Compute rank scores,
    top-$k$ accuracy, Brier scores, and credible set sizes, both for GNN models
    and for non-ML baselines.
\end{enumerate}

\section{Research Hypotheses}

The experiments are designed to test three specific hypotheses:

\begin{description}
  \item[H1 -- Causal Path Integrity.] Models that preserve higher-order causal
    paths (DBGNN, DAG-GNN) will outperform models that operate on the static
    projection (StaticGNN), because temporal ordering strictly constrains which
    transmission paths are causally feasible.
  \item[H2 -- Granularity Convergence.] As the transmission probability $\beta$
    decreases, outbreaks become sparser and more dependent on the precise sequence
    of contacts. The performance gap between temporally-aware models and the
    static baseline therefore increases with decreasing $\beta$.
  \item[H3 -- Temporal Signal-to-Noise Resilience.] When observations are made
    early (small fraction of $T$ elapsed), models that exploit temporal structure
    degrade more gracefully than static models, because they can use partial
    causal chain information even when the infection has not yet spread widely.
\end{description}

% ============================================================
\chapter{The SIR Epidemic Model}
\label{chap:sir}
% ============================================================

\section{Compartmental Models}

The SIR model partitions the node population at every time step into three
mutually exclusive compartments:

\begin{itemize}[noitemsep]
  \item $\mathbf{S}$ (Susceptible): healthy nodes that can acquire the infection.
  \item $\mathbf{I}$ (Infectious): infected nodes that can transmit to susceptible
    neighbours.
  \item $\mathbf{R}$ (Recovered / Removed): nodes that have recovered and are
    permanently immune.
\end{itemize}

On a network the mean-field ordinary differential equations are replaced by
stochastic individual-level transition rules. On temporal contact sequences the
transitions are further conditioned on whether the relevant edge is active at
the time of each potential event.

\section{Continuous-Time SIR on Static Networks}

On a static network with adjacency matrix $\mathbf{A}$, the continuous-time SIR
process assigns two independent Poisson processes to each pair of nodes:

\begin{align}
  S_v \xrightarrow{\beta \cdot \mathbf{A}_{uv}} I_v & \quad \text{(infection from neighbour } u \text{)} \\
  I_v \xrightarrow{\mu} R_v & \quad \text{(spontaneous recovery)}
\end{align}

Each $S$-$I$ contact edge fires at rate $\beta$; each $I$ node recovers at rate
$\mu$. The basic reproduction number for a homogeneous random network is:
\begin{equation}
  R_0 = \frac{\beta}{\mu} \langle k \rangle,
  \label{eq:r0}
\end{equation}
where $\langle k \rangle$ is the mean degree. An outbreak is sustained (endemic)
when $R_0 > 1$.

\section{Discrete-Time SIR on Temporal Contact Sequences}

The simulation engine in \texttt{tsir/} uses the discrete-time, event-driven
formulation of Holme (2020). At each contact event $(u, v, t)$ drawn from the
temporal sequence:

\begin{enumerate}
  \item If $u \in I$ and $v \in S$: node $v$ is infected with probability $\beta$
    (Bernoulli trial).
  \item If $v \in I$ and $u \in S$: node $u$ is infected with probability $\beta$
    (contacts are undirected).
  \item After processing all contacts at time $t$: each $I$ node recovers with
    probability $\mu$ (a single Bernoulli trial per infectious node per time step).
\end{enumerate}

\begin{insight}{Why Bernoulli, not Poisson?}
In continuous time infection and recovery are Poisson processes. The discrete
approximation replaces each infinitesimal time interval $dt$ with a single time
step. The per-contact transmission probability $\beta$ is therefore the
probability that a transmission event occurs at one specific contact, not the
rate of a Poisson process. The two are equivalent in the limit $dt \to 0$. The
discrete formulation is exact for event-driven temporal contact sequences where
each contact is an instantaneous event.
\end{insight}

\subsection{The Temporal Constraint}

The critical difference from a static network model is that node $v$ can only
be infected by node $u$ at a contact event $(u,v,t)$ if $u$ is currently in
the $I$ compartment at time $t$. This means:

\begin{itemize}[noitemsep]
  \item A contact $(u,v,t_1)$ where $u \in S$ at $t_1$ is harmless even if $u$
    later becomes infected.
  \item A contact $(u,v,t_2)$ where $u \in I$ at $t_2$ can transmit, but only
    if $v \in S$ at $t_2$ — it cannot retroactively infect $v$ for contacts at
    $t < t_2$.
\end{itemize}

This temporal constraint is precisely what the causally-aware graph
representations (De Bruijn graph, temporal event DAG) encode explicitly.

\section{SIR Parameters Used in Experiments}

\begin{table}[H]
\centering
\caption{Default SIR simulation parameters for the \texttt{france\_office}
  network. Other networks use identical parameters unless noted.}
\label{tab:sir_params}
\begin{tabular}{lll}
\toprule
Parameter & Symbol & Value \\
\midrule
Transmission probability & $\beta$ & 0.30 \\
Recovery probability & $\mu$ & 0.20 \\
Ground-truth runs per source & $n_{\text{runs}}$ & 5000 \\
Monte Carlo training runs per source & $n_{\text{mc}}$ & 500 \\
Observation window start & $t_{\text{start}}$ & 0 \\
Observation window end & $t_{\text{end}}$ & 274 \\
\bottomrule
\end{tabular}
\end{table}

The $\beta$-sweep experiment varies $\beta \in \{0.10, 0.20, 0.30, 0.50,
0.70\}$ while keeping $\mu = 0.20$ fixed, altering the effective reproduction
number $R_0$ accordingly.

\section{Source Filtering: Possible Sources}

Not every simulation run produces a useful training sample. Two binary masks
are computed:

\begin{description}
  \item[\texttt{maximal\_outbreak}{$[s,v]$}:] 1 if node $v$ became infected
    (i.e., is in $I \cup R$) in the \emph{maximal} (deterministic) outbreak
    starting from source $s$. A maximal outbreak ignores stochasticity and
    follows every possible contact edge — it is an upper bound on reachability.
  \item[\texttt{possible}{$[s,r,v]$}:] 1 if node $v$ is a \emph{possible
    source} for the observed state in run $r$. Formally, $v$ is a possible
    source if the observed infected set $I \cup R$ is contained in the reachable
    set from $v$ under the contact sequence.
\end{description}

The \texttt{lik\_possible} array encodes a log-likelihood correction:
\begin{equation}
  \texttt{lik\_possible}[s, r, v] =
  \begin{cases} 0 & \text{if } \texttt{possible}[s,r,v] = 1 \\ +\infty & \text{otherwise.} \end{cases}
\end{equation}
This term is subtracted from the model's log-probability scores during
inference, driving the score of topologically impossible sources to $-\infty$.

% ============================================================
\chapter{Temporal Contact Networks}
\label{chap:temporal_networks}
% ============================================================

\section{Formal Definition}

\begin{definition}{Temporal Contact Network}
A temporal contact network is a tuple $G = (\calV, \calE^t)$ where $\calV =
\{0,\ldots,N-1\}$ is a set of nodes and $\calE^t = \{(u,v,t) : u,v \in \calV,
t \in \mathbb{Z}_{\geq 0}\}$ is a multi-set of timestamped contact events. Each
event $(u,v,t)$ records that nodes $u$ and $v$ were in physical (or digital)
proximity at discrete time step $t$.
\end{definition}

In the implementation, temporal networks are stored as NetworkX graphs where
each edge $(u,v)$ carries a \texttt{times} list attribute containing all time
steps at which that pair was in contact:
\begin{lstlisting}[caption={Edge data structure for temporal graphs.}]
# H.edges(data=True) yields:
# (u, v, {'times': [t1, t2, t3, ...]})
for u, v, data in H.edges(data=True):
    contact_times = data['times']  # list of int
\end{lstlisting}

\section{Causal Walks}

\begin{definition}{Time-Respecting Path (Causal Walk)}
A time-respecting path (also called a causal walk) from node $s$ to node $v$
is a sequence of contact events $(v_0, v_1, t_1), (v_1, v_2, t_2), \ldots,
(v_{k-1}, v_k, t_k)$ such that $v_0 = s$, $v_k = v$, and $t_1 \leq t_2 \leq
\cdots \leq t_k$ (strictly: $t_i < t_{i+1}$ for transmission to be possible).
\end{definition}

Causal walks define which infection transmission paths are physically possible.
The static projection of a temporal network has far more paths than the causal
walk set, because the static graph ignores temporal ordering. The ratio of
causal walks to static paths grows smaller as the network becomes more temporal
(i.e., as the temporal ordering is more constraining). This motivates the
design of architectures that operate on the causal walk structure directly.

\section{Real-World Networks Used}

\begin{table}[H]
\centering
\caption{Networks used in experiments. All are empirical face-to-face contact
  networks from the SocioPatterns project unless noted otherwise.}
\label{tab:networks}
\begin{tabular}{llrrp{6cm}}
\toprule
Name & Key & $N$ & $|\calE^t|$ & Description \\
\midrule
toy\_holme & toy & $\sim$50 & $\sim$500 & Synthetic Holme-Kim network, used for
  fast debugging \\
karate\_static & karate & 34 & 78 & Zachary's Karate Club; static network used
  as a temporal network with all contacts at $t=0$ \\
france\_office & office & 232 & $\sim$9\,600 & InVS13 workplace contacts,
  274 hourly time steps \\
lyon\_ward & lyon & 75 & $\sim$32\,400 & Hospital ward contacts \\
malawi & malawi & 86 & $\sim$10\,000 & Rural village contacts, Malawi \\
students & students & 666 & $\sim$185\,000 & High-school student contacts \\
\bottomrule
\end{tabular}
\end{table}

% ============================================================
\chapter{Source Detection: Problem Definition and Likelihood}
\label{chap:source_detection}
% ============================================================

\section{Bayesian Formulation}

Applying Bayes' theorem to the source detection problem:
\begin{equation}
  P(\text{source} = s \mid \mathbf{O}) = \frac{P(\mathbf{O} \mid \text{source} = s) \, P(\text{source} = s)}{P(\mathbf{O})}.
  \label{eq:bayes}
\end{equation}
Under the uniform prior $P(\text{source} = s) = 1/N$ for all $s$, the
posterior is proportional to the likelihood $P(\mathbf{O} \mid \text{source} =
s)$. Computing this likelihood exactly requires summing over all stochastic
trajectories from source $s$ that are consistent with $\mathbf{O}$, which is
\#P-hard.

\section{Discriminative Surrogate}

Rather than computing the likelihood, we train a discriminative model:
\begin{equation}
  f_\theta : \mathbf{O} \longrightarrow \log \hat{P}_\theta(\text{source} = v \mid \mathbf{O}) \in \R^N,
\end{equation}
parameterised by $\theta$. The model is trained with the cross-entropy
(negative log-likelihood) loss:
\begin{equation}
  \calL(\theta) = -\frac{1}{|\mathcal{D}|} \sum_{(\mathbf{O}^{(i)}, s^{(i)}) \in \mathcal{D}} \log \hat{P}_\theta(\text{source} = s^{(i)} \mid \mathbf{O}^{(i)}),
  \label{eq:nll_loss}
\end{equation}
where $\mathcal{D}$ is the set of (observation, true source) pairs from the MC
simulations. This is identical to \texttt{F.nll\_loss} in PyTorch when the
model outputs log-probabilities.

\section{The Susceptible Masking Constraint}

A node $v$ that is observed in state $S$ (susceptible) at $t_{\text{end}}$
\emph{cannot} be the epidemic source, because the source by definition became
infectious at time $t=0$ and must therefore have left the $S$ compartment.
This hard constraint is enforced by:
\begin{equation}
  \text{score}(v) \leftarrow -\infty \quad \forall\, v : O_{v,S} = 1.
  \label{eq:susceptible_mask}
\end{equation}
All five models apply this mask before the final \texttt{log\_softmax},
guaranteeing zero probability for susceptible nodes. In code:
\begin{lstlisting}[caption={Susceptible masking (identical in all models).}]
# x: [B, N, 3], channel 0 = S state
susceptible_mask = x[..., 0].bool()   # True where node is S
scores = scores.masked_fill(susceptible_mask, float('-inf'))
return F.log_softmax(scores, dim=-1)  # [B, N]
\end{lstlisting}

\section{Log-Likelihood Correction}

After inference, the raw model log-probabilities are corrected by the
\texttt{lik\_possible} term:
\begin{equation}
  \tilde{\ell}(s \mid \mathbf{O}^{(r)}) = \log \hat{P}_\theta(s \mid \mathbf{O}^{(r)}) - \texttt{lik\_possible}[r, s].
  \label{eq:lik_correction}
\end{equation}
Since \texttt{lik\_possible} is $+\infty$ for topologically impossible sources
and $0$ for possible ones, this correction drives impossible sources to $-\infty$
and leaves possible sources unchanged. It implements an approximation to the
prior $P(\text{source} = s)$ based on network reachability.

% ============================================================
\chapter{Graph Neural Networks: Foundations}
\label{chap:gnn_foundations}
% ============================================================

\section{The Message Passing Framework}

All GNN architectures in this project are instances of the message passing
neural network (MPNN) framework of Gilmer et al. (2017). In the MPNN framework,
each layer $\ell$ updates node embeddings $\bh_v^{(\ell)}$ according to:

\begin{equation}
  \bh_v^{(\ell+1)} = \text{UPDATE}^{(\ell)}\!\left(\bh_v^{(\ell)},\;
    \bigoplus_{u \in \calN(v)} \text{MSG}^{(\ell)}\!\left(\bh_v^{(\ell)}, \bh_u^{(\ell)}, \mathbf{e}_{uv}\right)\right),
  \label{eq:mpnn}
\end{equation}
where:
\begin{itemize}[noitemsep]
  \item $\calN(v)$ is the set of neighbours of node $v$,
  \item $\mathbf{e}_{uv}$ is an optional edge feature vector,
  \item $\bigoplus$ is a permutation-invariant aggregation operator (sum, mean, or max),
  \item MSG and UPDATE are learnable functions (typically MLPs or linear layers).
\end{itemize}

After $L$ layers, the embedding $\bh_v^{(L)}$ aggregates information from the
$L$-hop neighbourhood of $v$.

\section{Specific Convolution Variants}

\subsection{GraphSAGE (SAGEConv)}

The mean-aggregator variant of GraphSAGE (Hamilton et al., 2017) is:
\begin{equation}
  \bh_v^{(\ell+1)} = \bW_1^{(\ell)} \bh_v^{(\ell)} + \bW_2^{(\ell)} \cdot \frac{1}{|\calN(v)|} \sum_{u \in \calN(v)} \bh_u^{(\ell)}.
  \label{eq:sage}
\end{equation}
SAGEConv is used in StaticGNN (optionally), TemporalGNN, DBGNN, and DAG-GNN.

\subsection{GraphConv}

The GraphConv layer (Morris et al., 2019), as implemented in PyTorch Geometric:
\begin{equation}
  \bh_v^{(\ell+1)} = \bW_1^{(\ell)} \bh_v^{(\ell)} + \bW_2^{(\ell)} \sum_{u \in \calN(v)} \bh_u^{(\ell)}.
  \label{eq:graphconv}
\end{equation}
This differs from SAGEConv only in the absence of mean-normalisation on the
neighbour sum, making it sensitive to node degree. StaticGNN defaults to
GraphConv with sum aggregation (\texttt{aggr: sum}).

\subsection{Why Output Log-Probabilities?}

All models output $\log \hat{P}(s \mid \mathbf{O}) \in \R^N$ (log-softmax).
This has two advantages:
\begin{enumerate}[noitemsep]
  \item The cross-entropy training loss (\texttt{F.nll\_loss}) expects
    log-probabilities and computes $-\log \hat{P}$ directly without an
    additional \texttt{log} call, improving numerical stability.
  \item Multiplying probabilities over independent runs (unlikely to be
    needed here, but standard practice) is equivalent to adding their
    logarithms, avoiding floating-point underflow.
\end{enumerate}

% ============================================================
\chapter{Graph Representations for Temporal Networks}
\label{chap:representations}
% ============================================================

Each GNN architecture requires a specific graph representation of the temporal
network $G = (\calV, \calE^t)$. All five representations are built once per
network (not per sample) by functions in \texttt{gnn/graph\_builder.py} and
stored in a \texttt{dict} of pre-computed tensors.

\section{Representation 1: Static Projection}
\label{sec:static_rep}

The static projection collapses all temporal edges into a single undirected
graph:
\begin{equation}
  \calE_{\text{static}} = \{(u,v) : \exists\, t \text{ s.t. } (u,v,t) \in \calE^t\}.
\end{equation}

The optional edge weight $w_{uv}$ counts the number of temporal contacts between
$u$ and $v$, normalised to $[0,1]$:
\begin{equation}
  w_{uv} = \frac{|\{t : (u,v,t) \in \calE^t\}|}{\max_{(u',v') \in \calE_{\text{static}}} |\{t : (u',v',t) \in \calE^t\}|}.
\end{equation}

The output dict has keys \texttt{edge\_index} $\in \mathbb{Z}^{2 \times 2|\calE|}$
(both directions) and optionally \texttt{edge\_weight} $\in \R^{2|\calE|}$.

\textbf{Loss of information}: The static projection loses all information about
\emph{when} each contact occurred. Two networks with identical static
projections but different contact orderings are indistinguishable to a GNN
operating on this representation.

\section{Representation 2: Temporal Activation Patterns}
\label{sec:activation_rep}

For the BacktrackingNetwork, each edge $(u,v)$ in the static projection is
additionally annotated with a binary activation vector:
\begin{equation}
  \mathbf{a}_{uv} \in \{0,1\}^T, \quad (\mathbf{a}_{uv})_t = \mathbf{1}[(u,v,t) \in \calE^t],
\end{equation}
where $T = t_{\max} + 1$ and $t_{\max} = \max_{(u,v,t) \in \calE^t} t$.

This gives the model access to \emph{which time steps} each edge was active, but
not directly to the ordering relative to other edges. The edge attribute matrix
is $\mathbf{A} \in \{0,1\}^{2|\calE| \times T}$.

For the \texttt{france\_office} network with $T=275$, each edge carries a
$275$-dimensional binary vector.

\section{Representation 3: Time-Sliced Snapshots}
\label{sec:snapshot_rep}

The temporal contact sequence is partitioned into non-overlapping time windows
of width $\Delta t$ (\texttt{group\_by\_time}):
\begin{equation}
  \text{snapshot } s = \{(u,v) : (u,v,t) \in \calE^t,\; \lfloor t/\Delta t \rfloor = s\}.
\end{equation}

The output is a dict \texttt{edge\_indeces}$[s]$ mapping each snapshot index $s$
to its edge\_index tensor. The TemporalGNN iterates over these snapshots in
\emph{reverse} temporal order (from $s = S-1$ down to $s = 0$), applying a
separate SAGEConv weight matrix per snapshot.

\section{Representation 4: De Bruijn Graph}
\label{sec:db_rep}

The De Bruijn graph $\mathcal{B}$ encodes higher-order causal paths in the
temporal network. It was introduced for temporal network analysis by
Qarkaxhija et al. (2022).

\subsection{Construction}

\begin{definition}{De Bruijn Graph of a Temporal Network}
Given temporal contact sequence $\calE^t$, the De Bruijn graph $\calB =
(\calV_\calB, \calE_\calB)$ is defined as:
\begin{itemize}
  \item \textbf{Event nodes}: For each contact $(u,v,t) \in \calE^t$, add a
    node $(u,v,t)$ to $\calV_\calB$.
  \item \textbf{Sentinel nodes}: For each original node $n \in \calV$ and each
    time $t \in \{t_{\min}, t_{\max}\}$, add a sentinel node $(n,n,t)$.
  \item \textbf{Causal edges}: Add directed edge $(x,y,t_1) \to (y,z,t_2)$
    whenever $t_2 > t_1$ (the walk $x \to y \to z$ is time-respecting).
  \item \textbf{Sentinel connections}: Sentinel node $(n,n,t)$ connects to all
    event nodes that start or end at node $n$.
\end{itemize}
\end{definition}

The De Bruijn graph is a \emph{line graph of the temporal network}: each node in
$\calB$ corresponds to a temporal edge (contact) in $G$, and each edge in $\calB$
corresponds to a time-respecting path of length 2 in $G$.

\subsection{Why Sentinel Nodes?}

Sentinel nodes $(n,n,t)$ serve as \emph{readout anchors}. After message passing
on $\calB$, the embedding of the end-time sentinel $(n,n,t_{\max})$ aggregates
all causal walk information that \emph{terminates at} node $n$ by time
$t_{\max}$. This provides a natural per-node readout without requiring a
separate pooling step.

\subsection{Tensor Layout}

The output dict contains:
\begin{itemize}[noitemsep]
  \item \texttt{db\_node\_to\_original} $\in \mathbb{Z}^{|\calV_\calB| \times 2}$:
    maps De Bruijn node $i$ to its original node pair $(u_{\text{orig}},
    v_{\text{orig}})$. For event node $(u,v,t)$: entry is $(u,v)$. For sentinel
    $(n,n,t)$: entry is $(n,n)$.
  \item \texttt{sentinel\_end\_indices} $\in \mathbb{Z}^N$:
    \texttt{sentinel\_end\_indices}$[n]$ is the integer index of De Bruijn node
    $(n,n,t_{\max})$ in $\calV_\calB$.
  \item \texttt{is\_sentinel} $\in \{0,1\}^{|\calV_\calB|}$: boolean mask; True
    for sentinel nodes (where $u_{\text{orig}} = v_{\text{orig}}$).
  \item \texttt{edge\_index} $\in \mathbb{Z}^{2 \times |\calE_\calB|}$: directed
    De Bruijn edges.
\end{itemize}

\section{Representation 5: Temporal Event DAG}
\label{sec:dag_rep}

The temporal event graph (TEG) of Saramäki et al. (2019) maps each directed
contact event to a node and connects causally compatible events by directed edges.

\subsection{Construction}

\begin{definition}{Temporal Event Graph (TEG)}
Given $\calE^t$ (treating undirected edges as two directed events), the TEG
$\mathcal{T} = (\calV_\mathcal{T}, \calE_\mathcal{T})$ is:
\begin{itemize}
  \item \textbf{Event nodes}: For each directed event $(u \to v, t)$ (both
    directions for undirected contacts), add a node $e_{uvt}$ to $\calV_\mathcal{T}$.
  \item \textbf{Causal edges}: Add directed edge $e_{uvt_1} \to e_{vwt_2}$ if
    $t_2 > t_1$ (and $t_2 - t_1 \leq \delta t$ if a maximum gap $\delta t$ is
    specified). The events share intermediate node $v$, so the infection could
    travel $u \to v$ at $t_1$ and then $v \to w$ at $t_2$.
\end{itemize}
\end{definition}

Because $t_2 > t_1$ strictly, the TEG has no directed cycles: it is a DAG.
This is the property that DAG-GNN exploits for backward propagation.

\subsection{Tensor Layout}

\begin{itemize}[noitemsep]
  \item \texttt{dag\_edge\_index} $\in \mathbb{Z}^{2 \times |\calE_\mathcal{T}|}$:
    forward causal edges $e_1 \to e_2$ with $t(e_1) < t(e_2)$.
  \item \texttt{event\_src\_node} $\in \mathbb{Z}^{|\calV_\mathcal{T}|}$:
    departing node $u$ for each event $e_{uvt}$.
  \item \texttt{event\_to\_node} $\in \mathbb{Z}^{|\calV_\mathcal{T}|}$:
    arriving node $v$ for each event $e_{uvt}$.
  \item \texttt{event\_times} $\in \mathbb{Z}^{|\calV_\mathcal{T}|}$: time $t$
    of each event.
\end{itemize}

\subsection{Complexity}

The number of event nodes is $|\calV_\mathcal{T}| = 2|\calE^t|$ (two directed
events per undirected contact). The number of DAG edges can be $O(|\calV_\mathcal{T}|^2)$
in the worst case (all events at the same node), but is typically much smaller
for real contact sequences. The $\delta t$ parameter limits the causal horizon
and reduces edge density.

\begin{table}[H]
\centering
\caption{Summary of the five graph representations and their properties.}
\label{tab:representations}
\begin{tabular}{lllll}
\toprule
Representation & Used by & Node count & Temporal info & Edge features \\
\midrule
Static projection & StaticGNN & $N$ & No & Optional weights \\
Activation patterns & BacktrackingNetwork & $N$ & Partial & $\mathbf{a}_{uv} \in \{0,1\}^T$ \\
Snapshots & TemporalGNN & $N$ & Ordering & None \\
De Bruijn graph & DBGNN & $\leq 2|\calE^t| + N$ & Full (causal walks) & None \\
Temporal event DAG & DAG-GNN & $2|\calE^t|$ & Full (DAG) & None \\
\bottomrule
\end{tabular}
\end{table}

% ============================================================
\chapter{StaticGNN: Baseline Architecture}
\label{chap:static_gnn}
% ============================================================

\section{Motivation}

StaticGNN provides the static-projection baseline: it is the strongest possible
GNN that ignores temporal ordering. All temporal information (contact times) is
discarded; the model sees only the aggregated static graph plus the SIR state
snapshot $\mathbf{O}$. If temporally-aware models do not outperform StaticGNN,
it suggests that temporal ordering is not needed for source detection on the
tested networks.

\section{Architecture}

The StaticGNN is a configurable module with four stages (implemented in
\texttt{gnn/static\_gnn.py}):

\begin{enumerate}
  \item \textbf{Preprocessing MLP}: $L_{\text{pre}}$ layers mapping node
    features $\bx_v \in \R^3$ to embeddings $\bh_v^{(0)} \in \R^{D_{\text{pre}}}$.
  \item \textbf{Graph Convolutions}: $L_{\text{conv}}$ GraphConv layers,
    each followed by optional batch normalization, PReLU activation, and dropout.
  \item \textbf{Postprocessing MLP}: $L_{\text{post}}$ fully-connected layers
    at dimension $D_{\text{hidden}}$ (identity in dimension).
  \item \textbf{Final Linear Layer}: Projects each node embedding to a scalar
    score, then reshapes to $[B, N]$ and applies log-softmax.
\end{enumerate}

\subsection{Input Representation}

The input to StaticGNN is a batch of SIR snapshots
$\mathbf{X} \in \R^{B \times N \times 3}$. Each node $v$ in batch item $b$ has
features:
\begin{equation}
  \bx_{b,v} = \bigl[O_{b,v,S},\; O_{b,v,I},\; O_{b,v,R}\bigr]^\top \in \{0,1\}^3.
\end{equation}
This is a one-hot encoding of the SIR state.

For PyTorch Geometric (PyG) style batching, the $[B, N, 3]$ tensor is
flattened to $[B \cdot N, 3]$ with a batch vector $\mathbf{b} \in \{0,\ldots,B-1\}^{B \cdot N}$
and the edge index is replicated $B$ times with appropriate offsets.

\subsection{Graph Convolution}

The default aggregation in the experiments is \texttt{aggr: sum} (GraphConv):
\begin{equation}
  \bh_v^{(\ell+1)} = \text{PReLU}\!\left(\text{BN}\!\left(\bW_1^{(\ell)} \bh_v^{(\ell)} + \bW_2^{(\ell)} \sum_{u \in \calN(v)} \bh_u^{(\ell)}\right)\right).
  \label{eq:static_gnn_conv}
\end{equation}
With batch normalization and PReLU, a dropout mask is applied after each layer
during training.

\subsection{Skip Connection}

When \texttt{skip: true} is set, the raw input features $\bx_v$ are
concatenated to the post-processed node embedding before the final linear layer:
\begin{equation}
  \text{score}(v) = \bW_{\text{out}} \cdot \bigl[\bh_v^{(L)},\, \bx_v\bigr]^\top.
\end{equation}
This allows the final layer to directly combine topological evidence (from $\bh_v^{(L)}$)
with the node's own observed state (from $\bx_v$). In practice this consistently
improves performance because the susceptible masking constraint means the node's
own state is highly informative.

\subsection{Output}

After applying the final layer to each node, the tensor is reshaped from
$[B \cdot N, 1]$ to $[B, N]$, the susceptible mask is applied, and
log-softmax is computed row-wise:
\begin{equation}
  \hat{\ell}_{b,v} = \log \hat{P}_\theta(\text{source} = v \mid \mathbf{O}_b).
\end{equation}

\section{Configuration}

\begin{lstlisting}[language=bash,caption={StaticGNN config for \texttt{france\_office} (from \texttt{exp/france\_office/static\_gnn.yml}).}]
model: static_gnn

train:
  n_mc: 500
  reps: 3
  test_size: 0.30
  batch_size: 128
  epochs: 500
  patience: 10
  lr: 0.001
  weight_decay: 0.0005
  seed: 42

static_gnn:
  num_preprocess_layers: 0
  embed_dim_preprocess: 16
  num_postprocess_layers: 0
  num_conv_layers: 4
  aggr: sum
  hidden_channels: 64
  dropout_rate: 0.2
  batch_norm: true
  skip: true
  use_edge_weights: false
\end{lstlisting}

\section{Complexity}

\begin{itemize}[noitemsep]
  \item \textbf{Parameters}: $O(D^2)$ per convolution layer, where $D = 64$ (hidden channels).
  \item \textbf{Forward pass time}: $O(B \cdot (N \cdot D^2 + |\calE| \cdot D))$ per layer.
  \item \textbf{Memory}: $O(B \cdot N \cdot D)$ for node embeddings, $O(|\calE|)$ for
    the static edge index (shared across all batch items).
\end{itemize}

% ============================================================
\chapter{TemporalGNN: Time-Slice Approach}
\label{chap:temporal_gnn}
% ============================================================

\section{Motivation}

TemporalGNN (implemented in \texttt{gnn/temporal\_gnn.py}) introduces temporal
information through \emph{snapshot processing}: it applies a separate SAGEConv
layer to each time-slice of the contact sequence. The key idea is that the
network topology at each time step is different, so information should propagate
along the graph structure that was actually present at that time.

\section{Architecture}

\begin{equation}
  \bh_v^{(t)} = \text{SAGEConv}_t\!\left(\bh_v^{(t-1)},\; \{(u,v) \in \text{snapshot } t\}\right),
  \quad t = S-1, S-2, \ldots, 0,
\end{equation}
where $S$ is the total number of snapshots and the convolutions are applied in
\emph{reverse} temporal order. Reverse ordering means that the model first
processes the contacts that occurred latest (close to $t_{\text{end}}$ when the
observation is made) and then propagates backward through earlier contacts. This
matches the direction of source detection: from the observed state backward
toward the probable source.

\subsection{Layer Structure}

\begin{enumerate}
  \item \textbf{Pre-linear}: $\bh_v^{(-1)} = \text{ReLU}(\bW_{\text{pre}} \bx_v)$,
    projecting from 3 to \texttt{hidden\_channels}.
  \item \textbf{Snapshot convolutions}: For each snapshot $t$ in reverse order,
    apply $\text{SAGEConv}_t$ (each has its own weight matrices $\bW_1^{(t)},
    \bW_2^{(t)}$).
  \item \textbf{Post-linear}: $\text{score}(v) = \bW_{\text{post}} \bh_v^{(0)}$,
    projecting to a scalar. Output shape $[B, N, 1]$ is squeezed to $[B, N]$.
  \item Log-softmax over nodes.
\end{enumerate}

\subsection{Implementation Note}

The model has \emph{one SAGEConv layer per snapshot}, so the number of parameters
scales with the number of snapshots. For the \texttt{france\_office} network with
274 time steps and \texttt{group\_by\_time}=1, this means 274 separate SAGEConv
layers. In practice, the \texttt{group\_by\_time} parameter is used to reduce
this to a manageable number of snapshots.

\begin{lstlisting}[caption={TemporalGNN forward pass (from \texttt{gnn/temporal\_gnn.py}).}]
def forward(self, x, edge_indeces, edge_attr=None):
    x = F.relu(self.lin_pre(x))
    count = 0
    for t in reversed(edge_indeces.keys()):
        x = F.relu(self.convs[count](x, edge_indeces[t], edge_attr))
        count += 1
    x = self.lin_post(x)
    return F.log_softmax(x, dim=1)
\end{lstlisting}

\subsection{Key Design Decision: Shared vs. Separate Weights}

TemporalGNN uses \emph{separate} weight matrices per snapshot
(\texttt{ModuleList} of SAGEConv layers). An alternative would be to share
weights across snapshots (a single SAGEConv applied at each step). The
separate-weight design allows the model to learn snapshot-specific aggregation
patterns — for example, that contacts early in the epidemic are more informative
about the source than late contacts.

\section{Configuration}

\begin{lstlisting}[language=bash,caption={TemporalGNN config (\texttt{exp/france\_office/temporal\_gnn.yml}).}]
model: temporal_gnn

temporal_gnn:
  hidden_channels: 64
  group_by_time: 1
\end{lstlisting}

% ============================================================
\chapter{BacktrackingNetwork: Temporal Activation Kernels}
\label{chap:backtracking}
% ============================================================

\section{Paper Reference}

Ru et al., ``Inferring Patient Zero on Temporal Networks via Graph Neural
Networks,'' AAAI 2023. The implementation in \texttt{gnn/backtracking\_network.py}
follows Figure 2 and Equations 3--6 of the paper.

\section{Motivation}

The BacktrackingNetwork (BN) enriches the static-graph message passing with
\emph{edge textures}: binary vectors $\mathbf{a}_{uv} \in \{0,1\}^T$ encoding
\emph{when} each edge was active. The intuition is that the temporal activation
pattern of an edge is informative about whether transmission along that edge is
plausible given the observed SIR state: if edge $(u,v)$ was only active at
$t < t_{\text{infection of } u}$, it cannot have transmitted the infection from
$u$ to $v$.

\section{Architecture}

\subsection{Initialisation}

\begin{align}
  \bh_v^{(0)} &= p^v(\bx_v) = \text{ReLU}(\bW_v \bx_v), \quad \bW_v \in \R^{D \times 3},
  \label{eq:bn_init_node} \\
  \bg_{(u,v)}^{(0)} &= p^e(\mathbf{a}_{uv}) = \text{ReLU}(\bW_e \mathbf{a}_{uv}), \quad \bW_e \in \R^{D \times T}.
  \label{eq:bn_init_edge}
\end{align}

Node features $\bx_v \in \{0,1\}^3$ are projected to a $D$-dimensional hidden
space. Edge activation patterns $\mathbf{a}_{uv} \in \{0,1\}^T$ are projected
to the same hidden dimension.

Note that the edge projections $\bg_{(u,v)}^{(0)}$ are \emph{shared across the
batch} at initialisation, because the activation patterns are network properties
(not sample-specific). They become sample-specific at layer 1 when the node
features $\bh_v^{(0)}$ (which \emph{are} sample-specific) enter the edge update.

\subsection{Layer Update (Equations 3 and 4)}

For layer $\ell = 0, \ldots, L-1$:

\textbf{Edge update} (Eq.~3 of the paper):
\begin{equation}
  \bg_{(u,v)}^{(\ell+1)} = f^e\!\left(\bigl[\bg_{(u,v)}^{(\ell)},\, \bh_u^{(\ell)}\bigr]\right),
  \label{eq:bn_edge_update}
\end{equation}
where $f^e: \R^{2D} \to \R^D$ is a fully-connected layer with ReLU. The edge
embedding incorporates the current embedding of its \emph{source} node.

\textbf{Node update} (Eq.~4 of the paper):
\begin{equation}
  \bh_v^{(\ell+1)} = \text{ReLU}\!\left(f^v\!\left(\bh_v^{(\ell)}\right) + \sum_{u \in \calN(v)} \bg_{(u,v)}^{(\ell+1)}\right),
  \label{eq:bn_node_update}
\end{equation}
where $f^v: \R^D \to \R^D$ is a separate fully-connected layer with ReLU. The
node embedding accumulates incoming edge messages via sum aggregation.

\subsection{Why Sum Aggregation?}

The paper uses sum (not mean) aggregation. This makes the node embedding
sensitive to the degree of the node: a node with many incoming edges has a
larger aggregation term. For source detection this is appropriate because the
source typically has more infected neighbours (higher effective out-degree in
the infection tree).

\subsection{Output}

\begin{equation}
  \text{score}(v) = \bW_{\text{final}} \bh_v^{(L)} \in \R,
\end{equation}
then susceptible masking and log-softmax produce $\hat{\ell} \in \R^{B \times N}$.

\section{Complete Forward Pass}

\begin{lstlisting}[caption={BacktrackingNetwork forward pass (from \texttt{gnn/backtracking\_network.py}).}]
def forward(self, x, edge_index, edge_attr):
    # x: [B, N, 3], edge_attr: [E, T]
    B, N, _ = x.shape

    # 1. Initial projections
    h = self.p_v(x)                                 # [B, N, D]
    g = self.p_e(edge_attr)                         # [E,  D]
    g = g.unsqueeze(0).expand(B, -1, -1).contiguous()  # [B, E, D]

    # 2. L layers of BN convolution
    for conv in self.convs:
        h, g = conv(h, g, edge_index)

    # 3. Final projection to scalar score per node
    scores = self.final(h).squeeze(-1)              # [B, N]

    # 4. Expert knowledge: susceptible nodes -> -inf
    susceptible_mask = x[..., 0].bool()
    scores = scores.masked_fill(susceptible_mask, float('-inf'))
    return F.log_softmax(scores, dim=-1)            # [B, N]
\end{lstlisting}

\section{Configuration}

\begin{lstlisting}[language=bash,caption={BacktrackingNetwork config (\texttt{exp/france\_office/backtracking.yml}).}]
model: backtracking

train:
  n_mc: 500
  reps: 3
  test_size: 0.20
  batch_size: 128
  epochs: 500
  patience: 30
  lr: 0.001
  weight_decay: 0.0001
  seed: 42

backtracking:
  hidden_dim: 32
  num_layers: 6
\end{lstlisting}

\section{Memory and Complexity}

\begin{itemize}[noitemsep]
  \item \textbf{Edge projection}: $\bW_e \in \R^{D \times T}$ has $D \cdot T$
    parameters. For $D=32$, $T=275$: approximately 8\,800 parameters just for
    the edge projection.
  \item \textbf{Per-layer}: Each BNConvLayer has $f^e \in \R^{2D \times D}$ and
    $f^v \in \R^{D \times D}$, so $3D^2$ parameters. With $L=6$ layers: $18D^2$.
  \item \textbf{Memory}: The edge hidden tensor $\bg \in \R^{B \times E \times D}$
    grows linearly in batch size $B$ and edge count $E$. For large temporal
    networks (many edges), this can dominate memory usage.
\end{itemize}

% ============================================================
\chapter{DBGNN: De Bruijn Graph Neural Networks}
\label{chap:dbgnn}
% ============================================================

\section{Paper Reference}

Qarkaxhija, Perri, and Scholtes, ``De Bruijn Goes Neural: Causality-Aware
Graph Neural Networks for Time Series Data on Dynamic Graphs,''
arXiv:2209.08311 (2022). Our implementation in \texttt{gnn/dbgnn.py} adapts
this architecture for epidemic source detection.

\section{Core Idea: Higher-Order Causal Paths}

The central observation motivating DBGNN is that standard GNNs on the static
projection aggregate information from the \emph{topological} $k$-hop
neighbourhood, whereas for temporal networks the relevant neighbourhood is the
\emph{causal} $k$-hop neighbourhood — the set of nodes reachable by
time-respecting paths of length at most $k$.

By operating on the De Bruijn graph, whose nodes represent individual contact
events (temporal edges), DBGNN performs message passing along \emph{causal walks
of length 2}. A node $(x,y,t_1)$ in the De Bruijn graph can send messages to
$(y,z,t_2)$ if and only if $t_2 > t_1$, enforcing the causal constraint
directly in the graph structure.

\section{Feature Initialisation}

Each De Bruijn node $i$ with original node pair $(u_{\text{orig}},
v_{\text{orig}})$ is initialised differently depending on whether it is a
sentinel or an event node:

\textbf{Sentinel nodes} (where $u_{\text{orig}} = v_{\text{orig}} = n$):
\begin{equation}
  \bh_i^{(0)} = \text{ReLU}\!\left(\bW_{\text{sent}} \cdot \bx_{n}\right), \quad \bW_{\text{sent}} \in \R^{D \times 3}.
  \label{eq:db_init_sent}
\end{equation}
The projection uses only the SIR state of the single associated original node.

\textbf{Event nodes} (where $u_{\text{orig}} \neq v_{\text{orig}}$):
\begin{equation}
  \bh_i^{(0)} = \text{ReLU}\!\left(\bW_{\text{evt}} \cdot \bigl[\bx_{u_{\text{orig}}},\, \bx_{v_{\text{orig}}}\bigr]^\top\right), \quad \bW_{\text{evt}} \in \R^{D \times 6}.
  \label{eq:db_init_evt}
\end{equation}
The projection concatenates the SIR states of both endpoints of the contact
event, giving the model access to the states of both nodes that participated
in the contact.

In code, both projections are computed for all nodes, then blended using the
boolean mask \texttt{is\_sentinel}:
\begin{lstlisting}[caption={DBGNN feature initialisation.}]
x_u = x[:, u_orig, :]               # [B, db_N, 3]
x_v = x[:, v_orig, :]               # [B, db_N, 3]

h_sent  = F.relu(self.proj_sentinel(x_u))           # [B, db_N, D]
h_event = F.relu(self.proj_event(
    torch.cat([x_u, x_v], dim=-1)))                 # [B, db_N, D]

mask = is_sentinel.view(1, db_N, 1).expand(B, db_N, D)
h = torch.where(mask, h_sent, h_event)              # [B, db_N, D]
\end{lstlisting}

\section{DBGNNLayer: Message Passing on the De Bruijn Graph}

Each DBGNN layer implements a SAGEConv-style update directly on $[B, |\calV_\calB|, D]$
tensors, without PyTorch Geometric overhead:
\begin{equation}
  \bh_i^{(\ell+1)} = \text{ReLU}\!\left(\bW_{\text{self}}^{(\ell)} \bh_i^{(\ell)} + \bW_{\text{agg}}^{(\ell)} \cdot \frac{1}{\deg^-(i)} \sum_{j : j \to i} \bh_j^{(\ell)}\right),
  \label{eq:db_layer}
\end{equation}
where $\deg^-(i) = |\{j : (j,i) \in \calE_\calB\}|$ is the in-degree of De
Bruijn node $i$.

\subsection{Scatter-Add Implementation}

The mean aggregation is implemented via \texttt{scatter\_add\_} on a preallocated
buffer, avoiding any Python loops over edges:
\begin{lstlisting}[caption={DBGNNLayer forward (from \texttt{gnn/dbgnn.py}).}]
def forward(self, h, edge_index):  # h: [B, db_N, D]
    B, db_N, D = h.shape
    src, dst = edge_index[0], edge_index[1]

    h_src = h[:, src, :]                        # [B, E, D]
    dst_idx = dst.view(1, -1, 1).expand(B, -1, D)
    agg_sum = h.new_zeros(B, db_N, D)
    agg_sum.scatter_add_(1, dst_idx, h_src)

    deg = h.new_zeros(db_N)
    deg.scatter_add_(0, dst, torch.ones(E, device=h.device))
    deg = deg.clamp(min=1).view(1, db_N, 1)

    agg_mean = agg_sum / deg
    return F.relu(self.W_self(h) + self.W_agg(agg_mean))
\end{lstlisting}

\section{Readout}

After $L$ layers of message passing, the embedding of the \emph{end-time
sentinel} $(n, n, t_{\max})$ is used as the representation of original node $n$:
\begin{equation}
  \bh_n^{\text{node}} = \bh_{\text{sentinel\_end\_indices}[n]}^{(L)}.
  \label{eq:db_readout}
\end{equation}
This is a direct index lookup, not a pooling operation. The end-time sentinel
receives messages from all event nodes that terminate at node $n$ by time
$t_{\max}$, aggregating the full causal walk evidence for that node.

\begin{equation}
  \text{score}(n) = \bW_{\text{out}} \bh_n^{\text{node}} \in \R,
\end{equation}
then susceptible masking and log-softmax.

\section{Configuration}

\begin{lstlisting}[language=bash,caption={DBGNN config (\texttt{exp/france\_office/dbgnn.yml}).}]
model: dbgnn

train:
  n_mc: 500
  reps: 3
  test_size: 0.20
  batch_size: 64
  epochs: 500
  patience: 20
  lr: 0.001
  weight_decay: 0.0005
  seed: 42

dbgnn:
  hidden_channels: 64
  num_conv_layers: 3
  conv_type: sage
  dropout_rate: 0.2
\end{lstlisting}

\section{Complexity Discussion}

The De Bruijn graph has $|\calV_\calB| = O(|\calE^t|)$ nodes (one per contact
event plus sentinels). For the \texttt{france\_office} network with
$|\calE^t| \approx 9\,600$ unique temporal edges (counting both directions), this
gives approximately 19\,200 event nodes plus $N = 232$ sentinel nodes. The De
Bruijn edge count is bounded by $|\calE_\calB| \leq |\calV_\calB|^2$ in the
worst case but is typically $O(|\calV_\calB| \cdot \bar{d})$ where $\bar{d}$ is
the mean in-degree in $\calB$. The per-batch memory footprint is $O(B \cdot
|\calV_\calB| \cdot D)$, which can be large for high-resolution temporal networks.

% ============================================================
\chapter{DAG-GNN: Temporal Event Graph GNN}
\label{chap:dag_gnn}
% ============================================================

\section{Paper References}

\begin{enumerate}
  \item Rey, Shi, Buciluä, and Smola, ``Directed Acyclic Graph Convolutional
    Networks,'' arXiv:2506.12218 (2025).
  \item Saramäki and Holme, ``Weighted Temporal Event Graphs,''
    arXiv:1912.03904 (2019).
\end{enumerate}

The implementation in \texttt{gnn/dag\_gnn.py} adapts Rey et al.'s DAG
convolution to the temporal event graph representation of Saramäki et al.

\section{Motivation: Why a DAG?}

The temporal event graph (TEG) is naturally a DAG because causal edges point
strictly forward in time. This has important consequences for message passing:

\begin{enumerate}
  \item \textbf{No cycles}: Information cannot loop back through earlier events.
    This means we can process the graph in a single topological sweep.
  \item \textbf{Backward propagation}: By reversing the DAG edges, information
    flows from \emph{late events} (near $t_{\text{end}}$) toward \emph{early
    events} (near $t=0$, where the source is). This aligns the information flow
    with the source detection task.
  \item \textbf{Exact inference}: On a DAG, belief propagation is exact (no
    loops to iterate over). While our model is approximate (learned weights,
    fixed layers), the DAG structure removes one source of approximation error.
\end{enumerate}

\section{Feature Initialisation}

Each event node $e_{uvt}$ (representing the directed contact $u \to v$ at time
$t$) is initialised by concatenating the SIR states of the source and
destination nodes:
\begin{equation}
  \bh_e^{(0)} = \text{ReLU}\!\left(\bW_{\text{proj}} \cdot [\bx_u, \bx_v]^\top\right), \quad \bW_{\text{proj}} \in \R^{D \times 6}.
  \label{eq:dag_init}
\end{equation}

\begin{lstlisting}[caption={DAG-GNN event feature initialisation.}]
x_src = x[:, ev_src, :]              # [B, n_events, 3]
x_dst = x[:, ev_dst, :]              # [B, n_events, 3]
h = F.relu(self.proj_event(
    torch.cat([x_src, x_dst], dim=-1)))  # [B, n_events, D]
\end{lstlisting}

\section{DAG Reversal and Backward Propagation}

The forward causal edges are reversed before message passing:
\begin{equation}
  \text{rev\_edge\_index} = \text{flip}(\text{dag\_edge\_index},\, \text{dim}=0),
\end{equation}
so an edge $e_1 \to e_2$ in the forward DAG (meaning $t(e_1) < t(e_2)$) becomes
$e_2 \to e_1$ in the reversed DAG. In the reversed DAG, information flows from
late events to early events.

\section{DAGConvLayer}

\begin{equation}
  \bh_e^{(\ell+1)} = \text{ReLU}\!\left(\bW_{\text{self}}^{(\ell)} \bh_e^{(\ell)} + \bW_{\text{agg}}^{(\ell)} \cdot \frac{1}{\deg^-_{\text{rev}}(e)} \sum_{e' : e' \to e \text{ in rev}} \bh_{e'}^{(\ell)}\right).
  \label{eq:dag_layer}
\end{equation}

In the reversed DAG, the in-neighbours of event $e$ are all events $e'$ such
that $e \to e'$ in the \emph{forward} DAG (i.e., $e'$ happens after $e$ and is
causally enabled by $e$). Thus, the aggregation at $e$ collects evidence from
all events that $e$ could have caused.

The implementation is structurally identical to DBGNNLayer, operating on the
reversed edge index.

\section{Readout: Events to Original Nodes}

After $L$ DAGConv layers, event embeddings are aggregated back to original
nodes. For node $v$, we aggregate all events where node $v$ was the
\emph{arriving} node (i.e., $v$ received the contact):
\begin{equation}
  \bh_v^{\text{node}} = \frac{1}{|\{e : \text{event\_to\_node}[e] = v\}|} \sum_{e : \text{event\_to\_node}[e] = v} \bh_e^{(L)}.
  \label{eq:dag_readout}
\end{equation}

The rationale for aggregating over \emph{arriving} events (not departing) is
that the source node $s$ is the origin of transmission chains — all events in
which information from $s$ eventually propagates will be on paths leading
\emph{to} the infected nodes. Aggregating over arriving events at each node
captures the accumulated causal evidence pointing toward that node as a
potential source.

\begin{lstlisting}[caption={DAG-GNN readout: events to nodes.}]
node_emb = torch.zeros(B, N, D, device=x.device)
dst_idx  = ev_dst.view(1, -1, 1).expand(B, n_events, D)
node_emb.scatter_add_(1, dst_idx, h)        # sum over arriving events

# Mean normalisation
count = torch.zeros(N, device=x.device)
count.scatter_add_(0, ev_dst, torch.ones(n_events, device=x.device))
count = count.clamp(min=1).view(1, N, 1)
node_emb = node_emb / count
\end{lstlisting}

\section{Handling Isolated Nodes}

Nodes that have no arriving events (i.e., no temporal contacts pointed at them)
receive a zero embedding from the scatter-add. The \texttt{clamp(min=1)} in the
mean normalisation prevents division by zero. Their score after the final linear
layer will be zero before the susceptible mask is applied; if they are
susceptible they will be masked to $-\infty$, otherwise they will receive a
finite but relatively low score.

\section{Configuration}

\begin{lstlisting}[language=bash,caption={DAG-GNN config (\texttt{exp/france\_office/dag\_gnn.yml}).}]
model: dag_gnn

train:
  n_mc: 500
  reps: 3
  test_size: 0.20
  batch_size: 64
  epochs: 500
  patience: 20
  lr: 0.001
  weight_decay: 0.0005
  seed: 42

dag_gnn:
  hidden_channels: 64
  num_conv_layers: 3
  dropout_rate: 0.2
  agg: mean
  delta_t: null
\end{lstlisting}

\section{Effect of the $\delta t$ Parameter}

When \texttt{delta\_t} is set to a finite integer $\delta$, only causal edges
$(e_1, e_2)$ with $t(e_2) - t(e_1) \leq \delta$ are included in the DAG. This
limits the temporal reach of causal connections and can significantly reduce the
edge count for dense contact sequences. Setting $\delta t = 1$ restricts to
\emph{immediately consecutive} events, which is the most conservative
interpretation of causality in a discrete-time model.

% ============================================================
\chapter{Training Pipeline}
\label{chap:training}
% ============================================================

\section{Data Structures}

\subsection{TSIRData}

The \texttt{TSIRData} dataclass (defined in \texttt{setup/data\_loader.py}) holds
all simulation outputs for a single network and set of SIR parameters:

\begin{lstlisting}[caption={TSIRData dataclass fields.}]
@dataclass
class TSIRData:
    config: dict            # original YAML config dict
    n_nodes: int            # N
    n_runs: int             # ground-truth runs per source
    mc_runs: int            # MC training runs per source
    mc_S:    np.ndarray     # [N, mc_runs, N] int8 -- S state
    mc_I:    np.ndarray     # [N, mc_runs, N] int8 -- I state
    mc_R:    np.ndarray     # [N, mc_runs, N] int8 -- R state
    truth_S: np.ndarray     # [N, n_runs, N] int8 -- S state
    truth_I: np.ndarray     # [N, n_runs, N] int8 -- I state
    truth_R: np.ndarray     # [N, n_runs, N] int8 -- R state
    maximal_outbreak: np.ndarray  # [N, N]
    possible: np.ndarray    # [N, n_runs, N]
    lik_possible: np.ndarray      # [N, n_runs, N]
\end{lstlisting}

The index convention for all three-dimensional arrays is
\texttt{[source, run, node]}: entry \texttt{mc\_S[s, r, v]} equals 1 if node
$v$ is in state $S$ in run $r$ of the MC simulation started from source $s$.

\subsection{SIRDataset: Constructing Training Samples}

The MC simulations are unrolled into a flat dataset of (observation, label)
pairs:
\begin{lstlisting}[caption={Flattening MC simulations into training samples.}]
# mc_S: [N, mc_runs, N] -> reshape to [N*mc_runs, N]
n_total = n_nodes * mc_runs
S = mc_S.reshape(n_total, n_nodes)
I = mc_I.reshape(n_total, n_nodes)
R = mc_R.reshape(n_total, n_nodes)
X = np.stack([S, I, R], axis=-1)   # [N*mc_runs, N, 3]

# Source label: repeat each source index mc_runs times
y = np.repeat(np.arange(n_nodes), mc_runs)  # [N*mc_runs]
# y[s * mc_runs + r] = s (source index)
\end{lstlisting}

The dataset has $N \times n_{\text{mc}}$ samples. Each sample is a tuple
$(\mathbf{X}_i, y_i)$ where $\mathbf{X}_i \in \{0,1\}^{N \times 3}$ is the
SIR state snapshot and $y_i \in \{0,\ldots,N-1\}$ is the true source index.

\section{Train/Validation Split}

The dataset is split with a stratified train/test split (scikit-learn
\texttt{train\_test\_split}) with \texttt{test\_size=0.20} or \texttt{0.30}
depending on the model config. Stratification is by source node index $y$,
ensuring that each source is represented in both train and validation sets.

Index pairs $(s, r)$ are the unit of splitting, not individual samples,
preserving the dependency structure within a run.

\section{Training Loop}

\begin{algorithm}[H]
\caption{Training Loop with Early Stopping}
\SetAlgoLined
\KwIn{Dataset $(X, y)$, model $f_\theta$, epochs $E$, patience $P$, batch size $B$}
\KwOut{Best model parameters $\theta^*$}
Split $(X_{\text{train}}, y_{\text{train}}), (X_{\text{val}}, y_{\text{val}})$ from $(X, y)$\;
$\theta^* \leftarrow \theta$, $\mathcal{L}^*_{\text{val}} \leftarrow +\infty$, counter $\leftarrow 0$\;
\For{epoch $= 1, \ldots, E$}{
  Shuffle $X_{\text{train}}$\;
  \For{each mini-batch $\{(\mathbf{X}_i, y_i)\}_{i \in \mathcal{B}}$}{
    $\hat{\ell}_i = f_\theta(\mathbf{X}_i)$ \tcp*{forward pass}
    $\mathcal{L} = -\frac{1}{|\mathcal{B}|}\sum_i \hat{\ell}_i[y_i]$ \tcp*{NLL loss}
    $\theta \leftarrow \theta - \eta \nabla_\theta \mathcal{L}$ \tcp*{Adam step}
  }
  $\mathcal{L}_{\text{val}} = \frac{1}{|X_{\text{val}}|}\sum_i -\hat{\ell}_i[y_i]$\;
  \If{$\mathcal{L}_{\text{val}} < \mathcal{L}^*_{\text{val}}$}{
    $\theta^* \leftarrow \theta$, $\mathcal{L}^*_{\text{val}} \leftarrow \mathcal{L}_{\text{val}}$, counter $\leftarrow 0$\;
  }
  \Else{
    counter $\leftarrow$ counter $+ 1$\;
    \If{counter $\geq P$}{
      Restore $\theta^*$\; \textbf{break}\;
    }
  }
}
\Return $\theta^*$\;
\end{algorithm}

The Adam optimiser is used with weight decay (L2 regularisation):
\begin{equation}
  \theta_{t+1} = \theta_t - \eta \cdot \hat{m}_t / (\sqrt{\hat{v}_t} + \epsilon) - \lambda \theta_t,
\end{equation}
with $\eta = 0.001$, $\lambda = 0.0005$ (or $0.0001$ for BacktrackingNetwork).

\section{Inference on Ground-Truth Simulations}

After training, the model is evaluated on the ground-truth simulations. The
$N \times n_{\text{runs}}$ snapshots from \texttt{truth\_S/I/R} are stacked into
a tensor $\mathbf{X}_{\text{truth}} \in \{0,1\}^{N \cdot n_{\text{runs}} \times N \times 3}$
and passed through the model in batches of 64 to produce log-probabilities.

The log-probabilities are exponentiated to obtain probabilities and then
corrected by \texttt{lik\_possible} (see Equation~\ref{eq:lik_correction}).

\section{Repetitions}

Each experiment is run \texttt{reps: 3} times with different random seeds for
weight initialisation and data shuffling. Evaluation metrics are reported as
mean and standard deviation across repetitions.

% ============================================================
\chapter{Evaluation Metrics}
\label{chap:metrics}
% ============================================================

All metrics are computed in \texttt{eval/scores.py}. Let $\mathbf{p}^{(r)} \in
[0,1]^N$ be the predicted probability vector for run $r$, and let $s^{(r)}$ be
the true source for run $r$. Let $\text{rank}^{(r)}$ be the rank (1-indexed) of
the true source $s^{(r)}$ in the descending ordering of $\mathbf{p}^{(r)}$:
\begin{equation}
  \text{rank}^{(r)} = 1 + |\{v : p_v^{(r)} > p_{s^{(r)}}^{(r)}\}|.
\end{equation}

Let $\text{sel}$ be the set of valid runs (those with at least
\texttt{min\_outbreak} infected nodes, default 2).

\section{Top-$k$ Accuracy}

\begin{equation}
  \text{Top-}k = \frac{1}{|\text{sel}|} \sum_{r \in \text{sel}} \mathbf{1}[\text{rank}^{(r)} \leq k].
\end{equation}
Reported for $k \in \{1, 3, 5\}$. Top-1 is the fraction of runs where the model
identifies the correct source as the single most probable candidate.

\section{Inverse Rank Score}

\begin{equation}
  \text{InvRank}(\delta) = \frac{1}{|\text{sel}|} \sum_{r \in \text{sel}} \frac{1 + \delta}{\text{rank}^{(r)} + \delta},
\end{equation}
where $\delta \geq 0$ is an offset parameter (\texttt{inverse\_rank\_offset}).
With $\delta = 0$: perfect score (rank 1) gives 1.0, rank 2 gives 0.5, rank
$N$ gives $1/N$.

\section{Normalised Brier Score}

\begin{equation}
  \text{BS}_{\text{norm}} = \frac{1}{|\text{sel}|} \sum_{r \in \text{sel}} \frac{1}{2/N} \cdot \frac{1}{N} \sum_v (p_v^{(r)} - \mathbf{1}[v = s^{(r)}])^2.
\end{equation}
The normalisation by $2/N$ makes the score 1.0 for a uniform distribution and
0 for a perfect prediction. Smaller is better.

\section{Credible Set Coverage}

The $p$-credible set for run $r$ is the smallest set $C_r$ of nodes such that
$\sum_{v \in C_r} p_v^{(r)} \geq p$. The credible set score is:
\begin{equation}
  \text{CS}(p) = \frac{1}{|\text{sel}|} \sum_{r \in \text{sel}} \mathbf{1}[s^{(r)} \in C_r],
\end{equation}
measuring the fraction of runs where the true source is contained in the
$p$-credible set. A well-calibrated model with $p=0.90$ should achieve
$\text{CS}(0.90) \approx 0.90$.

\section{Normalised Entropy}

\begin{equation}
  H_{\text{norm}} = \frac{1}{|\text{sel}|} \sum_{r \in \text{sel}} \frac{-\sum_v p_v^{(r)} \log p_v^{(r)}}{\log N}.
\end{equation}
This measures the concentration of the predicted distribution. Uniform gives
$H_{\text{norm}} = 1.0$; a point mass gives $H_{\text{norm}} = 0$.

\section{Practical Reporting}

In the \texttt{wandb} logs, metrics are reported under keys such as:
\begin{itemize}[noitemsep]
  \item \texttt{eval/top\_1}, \texttt{eval/top\_3}, \texttt{eval/top\_5}
  \item \texttt{eval/rank\_score\_off0} (inverse rank with $\delta=0$)
  \item \texttt{eval/brier\_score}, \texttt{eval/entropy}
\end{itemize}

% ============================================================
\chapter{Baselines}
\label{chap:baselines}
% ============================================================

Non-ML baselines are evaluated in \texttt{main\_eval.py} on the same
\texttt{TSIRData} artefact as the GNN models, ensuring a fair comparison.

\section{Uniform}

\begin{equation}
  p_v^{\text{unif}} = \frac{1}{N} \quad \forall\, v \in \calV.
\end{equation}
This is the trivial baseline (random guessing without information). Expected
Top-1 accuracy is $1/N$.

\section{Degree Centrality}

\begin{equation}
  p_v^{\text{deg}} = \frac{\deg_{\text{inf}}(v)}{\sum_u \deg_{\text{inf}}(u)},
\end{equation}
where $\deg_{\text{inf}}(v)$ is the degree of node $v$ in the subgraph induced
by infected nodes ($I \cup R$). The rationale is that epidemic sources tend to
be highly connected; high-degree nodes in the infected subgraph are more likely
to have initiated the spread.

\section{Closeness Centrality}

\begin{equation}
  p_v^{\text{close}} \propto \frac{N_{\text{inf}} - 1}{\sum_{u \in I \cup R} d(v, u)},
\end{equation}
where $d(v, u)$ is the shortest-path distance on the infected subgraph. Sources
are expected to be central (short paths to all infected nodes).

\section{Jordan Center}

\begin{definition}{Jordan Center}
The Jordan center of the infected subgraph is the node minimising the
eccentricity (maximum distance to any infected node):
\begin{equation}
  v^* = \argmin_{v \in I \cup R} \max_{u \in I \cup R} d(v, u).
  \label{eq:jordan_center}
\end{equation}
\end{definition}

The Jordan center is the standard network-science heuristic for source
detection, inspired by the result of Shah \& Zaman (2011) that rumour centrality
is maximised at the centroid of the infected subgraph on regular trees. On
irregular or temporal networks it is an approximation.

\section{Betweenness Centrality}

\begin{equation}
  p_v^{\text{betw}} \propto \sum_{s \neq v \neq t} \frac{\sigma_{st}(v)}{\sigma_{st}},
\end{equation}
where $\sigma_{st}$ is the number of shortest paths from $s$ to $t$ on the
infected subgraph and $\sigma_{st}(v)$ is the number of those passing through
$v$. Nodes controlling many shortest paths may be structural ``gatekeepers''
that are more likely to have seeded the epidemic.

\begin{table}[H]
\centering
\caption{Summary of non-ML baselines.}
\label{tab:baselines}
\begin{tabular}{lll}
\toprule
Baseline & Score & Complexity \\
\midrule
Uniform & $1/N$ for all nodes & $O(1)$ \\
Degree & Degree in infected subgraph & $O(|I \cup R|)$ \\
Closeness & Closeness centrality & $O(|I \cup R|^2)$ \\
Betweenness & Betweenness centrality & $O(|I \cup R|^3)$ \\
Jordan center & Minimise eccentricity & $O(|I \cup R|^2)$ \\
\bottomrule
\end{tabular}
\end{table}

% ============================================================
\chapter{Experimental Setup and Hypotheses}
\label{chap:experiments}
% ============================================================

\section{Network and SIR Parameter Combinations}

Each network in Table~\ref{tab:networks} is paired with a YAML config directory
under \texttt{exp/}. For example, \texttt{exp/france\_office/} contains
\texttt{tsir.yml}, \texttt{static\_gnn.yml}, \texttt{backtracking.yml},
\texttt{temporal\_gnn.yml}, \texttt{dbgnn.yml}, \texttt{dag\_gnn.yml}, and
\texttt{eval.yml}.

\subsection{Standard Experiments}

For each network, each of the five GNN models and five non-ML baselines is
evaluated under the default SIR parameters. The \texttt{france\_office} network
is the primary benchmark network because it is large enough to be non-trivial
($N=232$) but small enough to run all models in reasonable time.

\subsection{Sweep: Varying Transmission Rate $\beta$ (H2)}

\begin{table}[H]
\centering
\caption{$\beta$-sweep experimental conditions (all on \texttt{france\_office},
  $\mu = 0.20$ fixed).}
\label{tab:beta_sweep}
\begin{tabular}{lll}
\toprule
Directory & $\beta$ & Effective $R_0$ (approx.) \\
\midrule
\texttt{sweep\_vary\_beta/beta010} & 0.10 & $\sim$0.75 \\
\texttt{sweep\_vary\_beta/beta020} & 0.20 & $\sim$1.50 \\
\texttt{sweep\_vary\_beta/beta030} & 0.30 & $\sim$2.25 \\
\texttt{sweep\_vary\_beta/beta050} & 0.50 & $\sim$3.75 \\
\texttt{sweep\_vary\_beta/beta070} & 0.70 & $\sim$5.25 \\
\bottomrule
\end{tabular}
\end{table}

At $\beta = 0.10$ with $\mu = 0.20$, the effective $R_0 < 1$ and most outbreaks
die out quickly (sparse, localised spread). At $\beta = 0.70$, nearly the entire
network is infected in every run, making the true source difficult to identify
from network structure alone.

\textbf{Expected outcome (H2)}: At low $\beta$, only a small subset of the
causal paths was activated. Preserving the ordering (DBGNN, DAG-GNN) should
give a larger advantage over StaticGNN than at high $\beta$, where the static
projection is sufficient to capture the widespread infection.

\subsection{Sweep: Varying Observation Time (H3)}

\begin{table}[H]
\centering
\caption{Observation time sweep (all on \texttt{france\_office}).}
\label{tab:obs_sweep}
\begin{tabular}{lll}
\toprule
Directory & $t_{\text{end}}$ (fraction of $T$) & Fraction of time elapsed \\
\midrule
\texttt{sweep\_vary\_observation/t25} & 68 & 25\% \\
\texttt{sweep\_vary\_observation/t50} & 137 & 50\% \\
\texttt{sweep\_vary\_observation/t75} & 205 & 75\% \\
\texttt{sweep\_vary\_observation/t100} & 274 & 100\% \\
\bottomrule
\end{tabular}
\end{table}

Early observation (25\% of $T$ elapsed) means the epidemic has just started;
the infected set is small and the causal chains are short. Late observation
(100\%) means the epidemic has run its full course; many nodes are in $R$ and
the infected subgraph is large.

\textbf{Expected outcome (H3)}: At 25\% observation, temporal models should
retain relatively high accuracy because the short causal chains are captured
well by the De Bruijn or DAG structure. StaticGNN may degrade more severely
because the infected subgraph is too sparse to provide centrality signals.

\section{Wandb Artefact Lineage}

Every run in the pipeline is tracked in the \texttt{source-detection} wandb
project. The data flow follows:

\begin{enumerate}
  \item \texttt{main\_tsir.py} produces an artefact (e.g.,
    \texttt{france\_office:v0}) containing binary SIR arrays and the network
    pickle.
  \item Each model training run (\texttt{main\_train.py}) loads this artefact,
    trains, logs metrics, and saves probability tensors to
    \texttt{data/<run\_id>/probs\_rep\{n\}.pt}.
  \item Baseline evaluation (\texttt{main\_eval.py}) loads the same artefact
    and logs its metrics alongside the model runs.
\end{enumerate}

All runs are tagged with their model name (e.g., \texttt{model:dbgnn}) and
network name (e.g., \texttt{france\_office}) for easy filtering in the wandb
dashboard.

% ============================================================
\chapter{Implementation Details and Code Architecture}
\label{chap:implementation}
% ============================================================

\section{Model Registry}

The model registry (\texttt{gnn/registry.py}) provides a unified interface for
instantiating and running all five models. Each model is registered as a
\texttt{ModelSpec} containing four callables:

\begin{lstlisting}[caption={ModelSpec dataclass (from \texttt{gnn/registry.py}).}]
@dataclass
class ModelSpec:
    cls:        Type[torch.nn.Module]  # model class
    forward_fn: Callable               # (model, x_batch, graph_data, device) -> [B,N]
    builder_fn: Callable               # (H, **kwargs) -> dict
    build_fn:   Callable               # (model_cfg, n_nodes, graph_data) -> model
\end{lstlisting}

The training loop accesses a model only through its \texttt{ModelSpec}, keeping
the training code agnostic to model-specific details:
\begin{lstlisting}[caption={Generic training call via registry.}]
spec = get_model_spec(cfg.model)        # e.g. "dbgnn"
graph_data = spec.builder_fn(H)         # build graph representation
model = spec.build_fn(cfg.dbgnn, n_nodes, graph_data)
# ...
for x_batch, y_batch in loader:
    log_probs = spec.forward_fn(model, x_batch, graph_data, device)
    loss = F.nll_loss(log_probs, y_batch)
\end{lstlisting}

\section{Configuration System}

YAML configuration files are loaded into a \texttt{Config} object
(\texttt{setup/read\_config.py}) supporting dot-notation access:
\begin{lstlisting}[language=bash,caption={Structure of a complete experiment config.}]
model: dbgnn

nwk:
  type: empirical
  name: france_office
  t_max: 274

sir:
  beta: 0.30
  mu: 0.20
  start_t: 0
  end_t: 274
  n_runs: 5000
  mc_runs: 500

eval:
  min_outbreak: 2
  top_k: [1, 3, 5]
  inverse_rank_offset: [0]
  n_truth: 1000

train:
  n_mc: 500
  reps: 3
  test_size: 0.20
  batch_size: 64
  epochs: 500
  patience: 20
  lr: 0.001
  weight_decay: 0.0005
  seed: 42

dbgnn:
  hidden_channels: 64
  num_conv_layers: 3
  conv_type: sage
  dropout_rate: 0.2
\end{lstlisting}

\section{Binary Data Storage}

SIR state arrays are stored as flat binary files with \texttt{int8} dtype to
minimise disk space. For the \texttt{france\_office} network with $N=232$,
$n_{\text{mc}}=500$:
\begin{itemize}[noitemsep]
  \item \texttt{monte\_carlo\_S.bin}: $232 \times 500 \times 232 = 26\,912\,000$ bytes $\approx 25.7$ MB.
  \item Same for \texttt{\_I} and \texttt{\_R}.
  \item Ground-truth files: $232 \times 5000 \times 232 \approx 257$ MB each.
\end{itemize}

Loading is done with \texttt{np.fromfile} and immediate reshape:
\begin{lstlisting}[caption={Binary array loading (from \texttt{setup/data\_loader.py}).}]
mc_S = np.fromfile(f"data/{run_id}/monte_carlo_S.bin",
                   dtype=np.int8).reshape(n_nodes, -1, n_nodes)
# Shape: [N, mc_runs, N]
\end{lstlisting}

\section{Batching Strategies}

The five models use two different batching strategies:

\textbf{PyG-style batching} (StaticGNN): Node features are flattened to
$[B \cdot N, F]$ and the edge index is replicated $B$ times with offsets $\{0,
N, 2N, \ldots, (B-1)N\}$. A batch vector $\mathbf{b} \in \{0,\ldots,B-1\}^{B \cdot N}$
is used for graph-level pooling operations.

\textbf{Dense batching} (BacktrackingNetwork, DBGNN, DAG-GNN, TemporalGNN):
Node/event features are kept in $[B, N, D]$ tensors. The graph structure
(\texttt{edge\_index}) is shared and applied identically to each batch item.
Aggregations use \texttt{scatter\_add\_} on the batch dimension. This avoids
the overhead of PyG's sparse batching and is more memory-efficient when the
batch size is moderate.

\section{Device Management}

All models are written to be device-agnostic: graph data tensors (edge indices,
activation patterns) are moved to the appropriate device inside \texttt{forward\_fn},
and the training loop passes a \texttt{device} argument. The standard device
selection is:
\begin{lstlisting}
device = torch.device("cuda" if torch.cuda.is_available() else "cpu")
\end{lstlisting}

\section{Reproducibility}

Each training run sets three random seeds:
\begin{itemize}[noitemsep]
  \item \texttt{torch.manual\_seed(cfg.train.seed)}: PyTorch CPU operations.
  \item \texttt{torch.cuda.manual\_seed\_all(cfg.train.seed)}: GPU operations.
  \item \texttt{random.seed(cfg.train.seed)}: Python random (used in data shuffling).
\end{itemize}
With \texttt{reps: 3}, three independent repetitions are run with seeds
\texttt{cfg.train.seed}, \texttt{cfg.train.seed + 1}, and
\texttt{cfg.train.seed + 2}.

\section{Directory Structure Summary}

\begin{codebox}{Codebase Layout}
\begin{verbatim}
project-thesis-agents/
+-- gnn/
|   +-- static_gnn.py        # StaticGNN model
|   +-- temporal_gnn.py      # TemporalGNN model
|   +-- backtracking_network.py  # BacktrackingNetwork (Ru et al.)
|   +-- dbgnn.py             # DBGNN (Qarkaxhija et al.)
|   +-- dag_gnn.py           # DAG-GNN (Rey et al.)
|   +-- graph_builder.py     # Graph representation builders
|   +-- registry.py          # ModelSpec and MODEL_REGISTRY
|   +-- training.py          # Training loop utilities
|   `-- predict.py           # Inference utilities
+-- setup/
|   +-- data_loader.py       # TSIRData dataclass + wandb loading
|   +-- read_config.py       # YAML -> Config object
|   `-- read_network.py      # NetworkX graph loading
+-- eval/
|   `-- scores.py            # rank_score, top_k_score, etc.
+-- tsir/                    # C SIR simulator + Python wrapper
+-- exp/                     # Experiment configs (YAML)
|   +-- france_office/
|   +-- sweep_vary_beta/
|   `-- sweep_vary_observation/
+-- nwk/                     # Network data files
+-- data/                    # Simulation outputs (binary arrays)
+-- main_tsir.py             # Step 1: Run simulations
+-- main_train.py            # Step 2: Train GNN models
`-- main_eval.py             # Step 3: Evaluate baselines
\end{verbatim}
\end{codebox}

% ============================================================
\chapter{Conclusion}
\label{chap:conclusion}
% ============================================================

\section{Summary of Technical Content}

This guide has covered the complete technical stack of the master's thesis
project. The following is a condensed summary of the key points.

\textbf{Problem}: Given a temporal contact network and a snapshot of an SIR
epidemic state at $t_{\text{end}}$, identify the most probable source node.
This is a posterior inference problem that we solve discriminatively using
supervised GNNs trained on Monte Carlo simulations.

\textbf{Data}: The C-based temporal SIR simulator (adapted from Holme 2020)
produces ground-truth and MC simulation runs for every possible source node.
Data is stored as binary \texttt{int8} arrays and versioned via wandb artefacts.

\textbf{Representations}: Five graph representations are derived from the
temporal contact sequence, ranging from the fully information-losing static
projection to the causally-exact temporal event DAG. The representations form
a hierarchy of temporal information: static projection $\subset$ activation
patterns $\subset$ snapshots $\subset$ De Bruijn graph $\approx$ temporal event
DAG.

\textbf{Models}: All five models (StaticGNN, TemporalGNN, BacktrackingNetwork,
DBGNN, DAG-GNN) output $\log P(\text{source} = v \mid \mathbf{O}) \in \R^N$,
apply a hard susceptible mask, and are trained with NLL loss. They differ in
which graph representation they consume and how they implement message passing.

\textbf{Training}: Adam optimiser, early stopping on validation NLL, stratified
train/val split, 3 independent repetitions per experiment.

\textbf{Evaluation}: Five metrics (Top-$k$, inverse rank, Brier score, credible
set coverage, entropy) computed on 1000 ground-truth test runs, compared against
five non-ML baselines (uniform, degree, closeness, betweenness, Jordan center).

\textbf{Experiments}: Standard per-network evaluation plus two parameter sweeps
testing hypotheses H2 (varying $\beta$) and H3 (varying observation time).

\section{Model Comparison at a Glance}

\begin{table}[H]
\centering
\caption{Comparison of the five GNN architectures.}
\label{tab:model_comparison}
\begin{tabular}{llllp{4cm}}
\toprule
Model & Representation & Temporal info & Node count & Distinguishing feature \\
\midrule
StaticGNN & Static graph & None & $N$ & Flexible MLP + GConv, skip connections \\
TemporalGNN & Snapshots & Ordering & $N$ & Separate conv per snapshot \\
BacktrackingNet & Activation patterns & Partial & $N$ & Edge textures $\mathbf{a}_{uv} \in \{0,1\}^T$ \\
DBGNN & De Bruijn graph & Full causal walks & $O(|\calE^t|)$ & Higher-order causal paths \\
DAG-GNN & Temporal event DAG & Full DAG & $2|\calE^t|$ & Backward propagation on DAG \\
\bottomrule
\end{tabular}
\end{table}

\section{Open Questions and Future Directions}

\begin{enumerate}
  \item \textbf{Scalability}: DBGNN and DAG-GNN both have $O(|\calE^t|)$ nodes
    in their expanded graphs. For the \texttt{students} network with
    $|\calE^t| \approx 185\,000$, this produces $\approx 370\,000$ event nodes,
    potentially exceeding GPU memory for reasonable batch sizes. Sampling
    strategies or sparse approximations may be needed.
  \item \textbf{Unknown observation time}: The current formulation assumes
    $t_{\text{end}}$ is known. In practice one may observe a snapshot without
    knowing how long the epidemic has been running. Marginalising over $t_{\text{end}}$
    is an open problem.
  \item \textbf{Partial observation}: Only a fraction of nodes may have known
    states (sparse observation). Pinto et al. (2012) address this for Gaussian
    infection time models; the extension to GNNs is unexplored.
  \item \textbf{Multiple sources}: The current setup assumes a single source.
    Dong et al. (2019) study multi-source detection; temporal extensions are
    open.
  \item \textbf{HYPA-DBGNN}: Von Pichowski et al. (2024) propose using
    Hypergeometric Graph Ensembles to identify statistically anomalous temporal
    patterns in the De Bruijn graph, which could serve as a better inductive
    bias for source detection than the vanilla DBGNN.
\end{enumerate}

% ============================================================
% BIBLIOGRAPHY (abbreviated; expand with full BibTeX entries)
% ============================================================
\begin{thebibliography}{99}

\bibitem{holme2020}
  P.~Holme,
  ``Fast and general simulation of SIR epidemics on temporal contact networks,''
  \emph{PLoS Computational Biology}, 2021.

\bibitem{qarkaxhija2022}
  L.~Qarkaxhija, V.~Perri, and I.~Scholtes,
  ``De Bruijn Goes Neural: Causality-Aware Graph Neural Networks for Time Series
  Data on Dynamic Graphs,''
  arXiv:2209.08311, 2022.

\bibitem{rey2025}
  M.~Rey, Y.~Shi, O.~Buciluä, and A.~Smola,
  ``Directed Acyclic Graph Convolutional Networks,''
  arXiv:2506.12218, 2025.

\bibitem{saramaki2019}
  J.~Saramäki and P.~Holme,
  ``Weighted Temporal Event Graphs,''
  arXiv:1912.03904, 2019.

\bibitem{ru2023}
  Y.~Ru, A.~Hayhoe, and A.~Ortega,
  ``Inferring Patient Zero on Temporal Networks via Graph Neural Networks,''
  \emph{AAAI 2023}.

\bibitem{shah2011}
  D.~Shah and T.~Zaman,
  ``Rumors in a Network: Who's the Culprit?,''
  \emph{IEEE Transactions on Information Theory}, 57(8), 2011.

\bibitem{pinto2012}
  P.~C.~Pinto, P.~Thiran, and M.~Vetterli,
  ``Locating the Source of Diffusion in Large-Scale Networks,''
  \emph{Physical Review Letters}, 109:068702, 2012.

\bibitem{dong2019}
  M.~Dong, T.~Zheng, N.~Duffield, and T.~He,
  ``Multiple Rumor Source Detection with Graph Convolutional Networks,''
  \emph{ACM CIKM}, 2019.

\bibitem{hamilton2017}
  W.~Hamilton, Z.~Ying, and J.~Leskovec,
  ``Inductive Representation Learning on Large Graphs (GraphSAGE),''
  \emph{NeurIPS}, 2017.

\bibitem{gilmer2017}
  J.~Gilmer, S.~S.~Schütt, F.~Becker, M.~Müller, and K.-R.~Müller,
  ``Neural Message Passing for Quantum Chemistry (MPNN),''
  \emph{ICML}, 2017.

\bibitem{sterchi2025}
  M.~Sterchi, P.~Brack, and R.~Hilfiker,
  ``GNNs for Source Detection: Review and Benchmark,'' 2025.

\bibitem{vonpichowski2024}
  A.~von~Pichowski, I.~Scholtes, and collaborators,
  ``Inference of Sequential Patterns for Higher-Order Analysis of Temporal Event
  Data,'' 2024.

\end{thebibliography}

\end{document}
